\section{Introduction}
\label{sec:introduction}
Household batteries can absorb surplus rooftop photovoltaic generation and supply electricity when demand or import prices increase. By changing when households import and export power, their operation also affects the distribution network. Battery scheduling must therefore consider both household expenditure and electrical operating conditions. Dynamic export connections illustrate this requirement: Energex describes export allowances of up to 10\,kW per phase and an ordinary minimum of 1.5\,kW, subject to emergency exceptions~\cite{energex}. A controller must respond to its allowance while accounting for the wider consequences of simultaneous charging and discharging.

Reinforcement learning offers a way to adapt battery operation to recurring household conditions. Federated learning allows households to exchange model information while retaining local training records~\cite{fedavg}, and has been applied to building battery scheduling and transfer across household conditions~\cite{sievers}. This architecture reduces the need to centralise operating records, but it also creates a route through which compromised participants can affect other controllers. A malicious participant can corrupt local observations or rewards, or replace the model change transmitted to the server. Honest households may then incur costs through a shared policy learned from those contributions.

Robust aggregation limits suspicious contributions using coordinate statistics, neighbour consistency or agreement with a reference update~\cite{yin,krum,fltrust}. In household energy systems, these mechanisms must distinguish malicious deviations from legitimate heterogeneity. Solar capacity, battery size, consumption timing and local learning progress can all produce different updates from honest households. Rejecting an unusual contribution may therefore suppress useful information as well as an attack. Assessing a defence requires examining economic damage together with the treatment of honest contributions.

We investigate these outcomes using FedGrid, which groups changes by solar capacity, filters within groups, clips update norms and weights group means by agreement with a reference learner trained on disjoint households. The final configuration filters original changes before clipping at twice the median received-update norm. Four transmitted-update attacks and two local training-data attacks are evaluated in a testbed that couples battery controllers with export enforcement and unbalanced power-flow simulation. The evaluation includes household bills, solar curtailment, battery throughput, voltage and transformer loading.

The analysis addresses three questions: how aggregation changes bill damage across poisoning mechanisms; whether the combined filtering and clipping revision improves the observed attack outcomes across household datasets and transfer settings; and what a rejection log establishes about participation in the shared update. The last question requires distinguishing filter admission from positive aggregation weight, because an admitted contribution can subsequently receive zero weight.

The contributions are a coupled learning and feeder-simulation workflow, a comparison of ten aggregation configurations across six attacks, and an analysis relating recorded admission to realised aggregation coefficients. Matched-seed results show a lower observed mean Normalized-attack damage after the FedGrid revision, alongside persistent vulnerability to reward poisoning and modest bill gains over rule-based control. An audit identifies evaluation-to-training copies introduced by ACT preprocessing. We retain these qualifications when interpreting the results: the experiments characterise the implemented configurations and do not establish general robustness or fair household influence.

The following sections review related work, define the physical and learning models, and describe the evaluation and configuration-selection history. Results distinguish energy-management outcomes, attack damage and recorded rejection. The discussion examines their operational implications, while the appendix provides detailed tables, sensitivity analyses and reproducibility audits.
